\section{Related Work}
\label{sec:related}

\subsection{UWB positioning and NLOS mitigation}
UWB indoor positioning has been studied extensively; a comprehensive
survey is given in~\cite{alarifi2016uwb}. Fully open systems remain
comparatively rare: ATLAS provides an open time-difference-of-arrival
localizer~\cite{tiemann2017atlas}, and the UTIL dataset~\cite{zhao2024util}
releases labelled ranging logs for benchmarking. A dominant error source
in cluttered environments is non-line-of-sight propagation, in which an
obstructed direct path forces the signal along a longer route and biases
the range positively. Kolakowski and Modelski~\cite{kolakowski2017nlos}
mitigate this by classifying each link as line-of-sight or NLOS from the
DW1000 first-path power and feeding the decision to an extended Kalman
filter. Rather than a hard classification, we derive a continuous weight
from the first-path-to-total power ratio measured within the same range
report, which both attenuates suspect ranges in the position solve and
inflates their variance in the temporal filter.

\subsection{Anchor self-calibration}
The anchor coordinates required by any multilateration system are
conventionally measured by hand or with external instrumentation.
Self-calibration methods that instead infer the anchor layout from radio
measurements are surveyed in~\cite{ridolfi2021selfcal}. Recent work
scales anchor calibration to large environments using Gaussian
processes, but relies on a LiDAR-inertial odometry trajectory as
input~\cite{yuan2025uwbgp}. In contrast, the self-survey used here
requires only the pairwise ranges the anchors can already measure among
themselves, and recovers the constellation by classical multidimensional
scaling~\cite{torgerson1952mds}.
